\documentclass[aspectratio=169,10pt,fontset=fandol]{ctexbeamer}
\usepackage{amsmath,amssymb,mathtools,bm,mathrsfs}
\usepackage{graphicx,booktabs,array,tabularx,multirow}
\usepackage{tikz,pgfplots,xparse,xstring,mdframed}
\usetikzlibrary{arrows.meta,calc,angles,positioning,shapes.geometric,graphs}
\pgfplotsset{compat=1.18}
\usetheme{fduclassroom}
\setbeamerfont{normal text}{size=\normalsize}
\setbeamerfont{title}{size=\large}
\setbeamerfont{subtitle}{size=\normalsize}
\setbeamerfont{frametitle}{size=\large}
\setbeamerfont{framesubtitle}{size=\small}
\setbeamerfont{block title}{size=\normalsize}
\setbeamerfont{block body}{size=\normalsize}
\setbeamerfont{section in toc}{size=\normalsize}
\setbeamerfont{section in head/foot}{size=\footnotesize}
\setbeamerfont{author in head/foot}{size=\footnotesize}
\setbeamerfont{title in head/foot}{size=\footnotesize}
\setbeamerfont{date in head/foot}{size=\footnotesize}
\setbeamertemplate{navigation symbols}{}
\setbeamertemplate{itemize item}{\normalsize\textbullet}
\setbeamertemplate{itemize subitem}{\small\textbullet}
\setbeamercolor{structure}{fg=blue!65!black}
\setbeamercolor{block title}{fg=blue!65!black,bg=blue!8}
\setbeamercolor{block body}{fg=black,bg=blue!3}
\setbeamertemplate{headline}{\leavevmode
\begin{beamercolorbox}[wd=\paperwidth,ht=3.6ex,dp=1.5ex,center]{palette secondary}
{\footnotesize\strut\insertsectionhead}
\end{beamercolorbox}}

% ctex fontset=fandol; system fonts are not redistributed.

\newenvironment{proofdetail}{\par\noindent\textbf{证明：}\quad}{\par}
\renewenvironment{solution}{\par\noindent\textbf{解：}\quad}{\par}
\makeatletter
\@ifundefined{definition}{\newtheorem{definition}{定义}[section]}{}
\@ifundefined{theorem}{\@ifundefined{c@definition}{\newtheorem{theorem}[theorem]{定理}}{\newtheorem{theorem}[definition]{定理}}}{}
\@ifundefined{example}{\@ifundefined{c@definition}{\newtheorem{example}[theorem]{例题}}{\newtheorem{example}[definition]{例题}}}{}
\@ifundefined{exercise}{\theoremstyle{definition}\@ifundefined{c@definition}{\newtheorem{exercise}[theorem]{练习}}{\newtheorem{exercise}[definition]{练习}}\theoremstyle{plain}}{}
\@ifundefined{exercise}{\theoremstyle{definition}\newtheorem{exercise}[theorem]{练习}\theoremstyle{plain}}{}
\makeatother
\makeatletter
\newcommand{\mathagentsectionlabelcontext}{%
\protected@edef\@currentlabel{\thesection}%
\edef\@currentHref{Outline\the\c@part.\the\c@section}%
}
\newcommand{\mathagentexerciselabelcontext}[1]{%
\hypertarget<1>{#1}{}%
\edef\@currentHref{#1}%
}
\makeatother
\title{lecture.tex}
\begin{document}
\section{数学内容}
% FRAME-ID F001
\typeout{MATH-AGENT-FRAME-BEGIN F001}
\begin{frame}[t,fragile]{目录}
\normalsize
{\color{blue}\normalsize 1. 数学内容\par}\vspace{0.22em}
{\color{gray!40}\normalsize 2. 数列极限的定义\par}\vspace{0.22em}
{\color{gray!40}\normalsize 3. 极限的运算\par}\vspace{0.22em}
\end{frame}
\typeout{MATH-AGENT-FRAME-END F001}
% FRAME-ID F002
\typeout{MATH-AGENT-FRAME-BEGIN F002}
\begin{frame}[t,fragile]{[NO \textbackslash{}title GIVEN] [NO \textbackslash{}a}
\normalsize
{\normalsize
% SOURCE-BEGIN U0001 0 55

\maketitle
\par\noindent 离线合成演示：不代表真实模型对本专题的教学验收。\par

% SOURCE-END U0001
}
\end{frame}
\typeout{MATH-AGENT-FRAME-END F002}
\section{数列极限的定义}
% FRAME-ID F003
\typeout{MATH-AGENT-FRAME-BEGIN F003}
\begin{frame}[t,fragile]{目录}
\normalsize
{\color{black}\normalsize 1. 数学内容\par}\vspace{0.22em}
{\color{blue}\normalsize 2. 数列极限的定义\par}\vspace{0.22em}
{\color{gray!40}\normalsize 3. 极限的运算\par}\vspace{0.22em}
\end{frame}
\typeout{MATH-AGENT-FRAME-END F003}
% FRAME-ID F004
\typeout{MATH-AGENT-FRAME-BEGIN F004}
\begin{frame}[t,fragile]{称实数列 a\_n 收敛于 a，如果对每个 ε>0}
\normalsize
{\normalsize
% PRESENTATION-SCAFFOLD original heading label binds to section
\mathagentsectionlabelcontext
% SOURCE-BEGIN U0002 0 31
\label{sec:S01}
% UNIT S01.U01

% SOURCE-END U0002
}
\begin{block}{定义}
{\normalsize
% SOURCE-BEGIN U0003 0 89
称实数列 $a_n$ 收敛于 $a$，如果对每个 $\varepsilon>0$，存在正整数 $N$，使得所有 $n>N$ 都满足 $|a_n-a|<\varepsilon$。

% SOURCE-END U0003
}
\end{block}
\end{frame}
\typeout{MATH-AGENT-FRAME-END F004}
% FRAME-ID F005
\typeout{MATH-AGENT-FRAME-BEGIN F005}
\begin{frame}[t,fragile]{收敛实数列的极限唯一。 若极限 a b，取 ε=}
\normalsize
\begin{block}{定理}
{\normalsize
% SOURCE-BEGIN U0004 0 12
收敛实数列的极限唯一。

% SOURCE-END U0004
}
\end{block}
\pause
{\normalsize
\textbf{证明：}
% SOURCE-BEGIN U0005 0 136
若极限 $a\ne b$，取 $\varepsilon=|a-b|/3>0$。由两次收敛定义，充分大的 $n$ 同时满足 $|a_n-a|<\varepsilon$ 和 $|a_n-b|<\varepsilon$。三角不等式给出 $|a-b|<2|a-b|/3$，矛盾。

% SOURCE-END U0005
}
\end{frame}
\typeout{MATH-AGENT-FRAME-END F005}
% FRAME-ID F006
\typeout{MATH-AGENT-FRAME-BEGIN F006}
\begin{frame}[t,fragile]{证明 a\_n=1/n 收敛于 0。}
\normalsize
\begin{block}{例题}
{\normalsize
% SOURCE-BEGIN U0006 0 22
证明 $a_n=1/n$ 收敛于 $0$。

% SOURCE-END U0006
}
\end{block}
\end{frame}
\typeout{MATH-AGENT-FRAME-END F006}
% FRAME-ID F007
\typeout{MATH-AGENT-FRAME-BEGIN F007}
\begin{frame}[t,fragile]{解：任取 ε>0，取正整数 N>1/ε。当 n>}
\normalsize
{\normalsize
% SOURCE-BEGIN U0007 0 137
\par\noindent\textbf{解：}任取 $\varepsilon>0$，取正整数 $N>1/\varepsilon$。当 $n>N$ 时，$|1/n-0|=1/n<1/N<\varepsilon$。故所求极限为 $0$。\par
% UNIT S01.E01

% SOURCE-END U0007
}
\end{frame}
\typeout{MATH-AGENT-FRAME-END F007}
% FRAME-ID F008
\typeout{MATH-AGENT-FRAME-BEGIN F008}
\begin{frame}[t,fragile]{按定义证明 3/n→0。}
\normalsize
{\normalsize
% PRESENTATION-SCAFFOLD exercise number and answer target
\setbeamertemplate{theorems}[numbered]
\mathagentexerciselabelcontext{mathagent.exercise.U0008.0}
% SOURCE-BEGIN U0008 0 140
\begin{exercise}
\label{mathagent:exercise:0cc0e7ee849eb730aacd59475ce492febad40593b04fd779ce32c70f94e1140f}
按定义证明 $3/n\to0$。
\end{exercise}
% SOURCE-END U0008
}
\end{frame}
\typeout{MATH-AGENT-FRAME-END F008}
% FRAME-ID F009
\typeout{MATH-AGENT-FRAME-BEGIN F009}
\begin{frame}[t,fragile]{解：取正整数 N>3/ε，则 n>N 时 3/n}
\normalsize
{\normalsize
% SOURCE-BEGIN U0009 0 80
\par\noindent\textbf{解：}取正整数 $N>3/\varepsilon$，则 $n>N$ 时 $3/n<\varepsilon$。\par

% SOURCE-END U0009
}
\end{frame}
\typeout{MATH-AGENT-FRAME-END F009}
\section{极限的运算}
% FRAME-ID F010
\typeout{MATH-AGENT-FRAME-BEGIN F010}
\begin{frame}[t,fragile]{目录}
\normalsize
{\color{black}\normalsize 1. 数学内容\par}\vspace{0.22em}
{\color{black}\normalsize 2. 数列极限的定义\par}\vspace{0.22em}
{\color{blue}\normalsize 3. 极限的运算\par}\vspace{0.22em}
\end{frame}
\typeout{MATH-AGENT-FRAME-END F010}
% FRAME-ID F011
\typeout{MATH-AGENT-FRAME-BEGIN F011}
\begin{frame}[t,fragile]{若 a\_n→ a，则 r\_n=a\_n-a 称为余}
\normalsize
{\normalsize
% PRESENTATION-SCAFFOLD original heading label binds to section
\mathagentsectionlabelcontext
% SOURCE-BEGIN U0010 0 31
\label{sec:S02}
% UNIT S02.U01

% SOURCE-END U0010
}
\begin{block}{定义}
{\normalsize
% SOURCE-BEGIN U0011 0 46
若 $a_n\to a$，则 $r_n=a_n-a$ 称为余项，并且 $r_n\to0$。

% SOURCE-END U0011
}
\end{block}
\end{frame}
\typeout{MATH-AGENT-FRAME-END F011}
% FRAME-ID F012
\typeout{MATH-AGENT-FRAME-BEGIN F012}
\begin{frame}[t,fragile]{若实数列 a\_n→ a 且 b\_n→ b，则 a}
\normalsize
\begin{block}{定理}
{\normalsize
% SOURCE-BEGIN U0012 0 49
若实数列 $a_n\to a$ 且 $b_n\to b$，则 $a_n+b_n\to a+b$。

% SOURCE-END U0012
}
\end{block}
\pause
{\normalsize
\textbf{证明：}
% SOURCE-BEGIN U0013 0 112
任取 $\varepsilon>0$，分别对两个收敛数列使用 $\varepsilon/2$。充分大的 $n$ 满足 $|(a_n+b_n)-(a+b)|\leq |a_n-a|+|b_n-b|<\varepsilon$。

% SOURCE-END U0013
}
\end{frame}
\typeout{MATH-AGENT-FRAME-END F012}
% FRAME-ID F013
\typeout{MATH-AGENT-FRAME-BEGIN F013}
\begin{frame}[t,fragile]{求 lim\_n→∞(2+1/n)。}
\normalsize
\begin{block}{例题}
{\normalsize
% SOURCE-BEGIN U0014 0 30
求 $\lim_{n\to\infty}(2+1/n)$。

% SOURCE-END U0014
}
\end{block}
\end{frame}
\typeout{MATH-AGENT-FRAME-END F013}
% FRAME-ID F014
\typeout{MATH-AGENT-FRAME-BEGIN F014}
\begin{frame}[t,fragile]{解：常数列 2 的极限为 2，且 1/n→0。由}
\normalsize
{\normalsize
% SOURCE-BEGIN U0015 0 89
\par\noindent\textbf{解：}常数列 $2$ 的极限为 $2$，且 $1/n\to0$。由加法法则，所求极限为 $2$。\par
% UNIT S02.E01

% SOURCE-END U0015
}
\end{frame}
\typeout{MATH-AGENT-FRAME-END F014}
% FRAME-ID F015
\typeout{MATH-AGENT-FRAME-BEGIN F015}
\begin{frame}[t,fragile]{按定义证明 5/n→0。}
\normalsize
{\normalsize
% PRESENTATION-SCAFFOLD exercise number and answer target
\setbeamertemplate{theorems}[numbered]
\mathagentexerciselabelcontext{mathagent.exercise.U0016.0}
% SOURCE-BEGIN U0016 0 140
\begin{exercise}
\label{mathagent:exercise:f708d5de6eeb2a0d1d46cddf76922686b9f273c43bd0798e15293c1e0afd54ab}
按定义证明 $5/n\to0$。
\end{exercise}
% SOURCE-END U0016
}
\end{frame}
\typeout{MATH-AGENT-FRAME-END F015}
% FRAME-ID F016
\typeout{MATH-AGENT-FRAME-BEGIN F016}
\begin{frame}[t,fragile]{解：取正整数 N>5/ε，则 n>N 时 5/n}
\normalsize
{\normalsize
% SOURCE-BEGIN U0017 0 80
\par\noindent\textbf{解：}取正整数 $N>5/\varepsilon$，则 $n>N$ 时 $5/n<\varepsilon$。\par

% SOURCE-END U0017
}
\end{frame}
\typeout{MATH-AGENT-FRAME-END F016}
\end{document}
